% ભાગ: first-order-logic
% પ્રકરણ: beyond
% વિભાગ: higher-order-logic

\documentclass[../../../include/open-logic-section]{subfiles}

\begin{document}

\olfileid{fol}{byd}{hol}

\olsection{ઉચ્ચ-ક્રમ તર્કશાસ્ત્ર}

પ્રથમ-ક્રમ તર્કશાસ્ત્રમાંથી દ્વિતીય-ક્રમ તર્કશાસ્ત્ર તરફ જવાથી વિધિવત્
ભાષાની અંદર પ્રથમ-ક્રમ !!{domain}ની વસ્તુઓના ગણોની વાત કરી શક્યા.
ત્યાં જ શા માટે અટકવું? દાખલા તરીકે, તૃતીય-ક્રમ તર્કશાસ્ત્ર આપણને
વસ્તુઓના ગણોના ગણો, અથવા કદાચ વસ્તુઓ અને વસ્તુઓના ગણો બંને ધરાવતા ગણો
સાથે વ્યવહાર કરવા દેવું જોઈએ. અને ચતુર્થ-ક્રમ તર્કશાસ્ત્ર આપણને એ
પ્રકારની વસ્તુઓના ગણોની વાત કરવા દેશે. તમે ધાર્યું હશે તેમ, આ વિચારને
ઇચ્છા મુજબ પુનરાવર્તિત કરી શકાય છે.

વ્યવહારમાં ઉચ્ચ-ક્રમ તર્કશાસ્ત્ર ઘણી વાર સંબંધોને બદલે વિધેયોના પદોમાં
ઘડવામાં આવે છે. (કુદરતી એકરૂપીકરણોને બાદ કરીએ તો આ
તફાવત અગત્યનો નથી.) કેટલાક પાયાના ``વર્ગો'' $A$, $B$, $C$,~\dots
આપેલા હોય (જેને હવે આપણે ``પ્રકારો'' કહીશું), તો નીચેની શરત મૂકીને
નવા પ્રકારો બનાવી શકીએ:
\begin{quote}
જો $\sigma$ અને~$\tau$ સાન્ત પ્રકારો હોય, તો $\sigma \to \tau$ પણ
સાન્ત પ્રકાર છે.
\end{quote}
\sourcecorrection{OLBYD-006}{મૂળની
\texttt{higher-order-logic.tex}, પંક્તિ~21માં અંગ્રેજી ક્રિયાપદ
``formulated''ની અંદર નામપદ માટેનું શબ્દકોશ-ટોકન
\texttt{\string!\string!\{formula\}} મૂકીને
\texttt{\string!\string!\{formula\}ted} રચાયું છે.
અહીં ક્રિયાપદનો સુસંગત ગુજરાતી અનુવાદ ટોકન વિના આપ્યો છે.}
પ્રકારોને એવી વાક્યરચનાત્મક ``નામપટ્ટીઓ'' તરીકે વિચારો જે આપણા
!!{domain}માં જોઈતી વસ્તુઓનું વર્ગીકરણ કરે છે; $\sigma \to \tau$ એવી
વસ્તુઓનું વર્ણન કરે છે જે પ્રકાર~$\sigma$ની વસ્તુઓને પ્રકાર~$\tau$ની
વસ્તુઓમાં મોકલતા વિધેયો છે. દાખલા તરીકે, ``સત્ય'' અને ``અસત્ય'' એવા
સત્યમૂલ્યોનો પ્રકાર~$\Omega$ અને પ્રાકૃતિક સંખ્યાઓનો પ્રકાર~$\Nat$
રાખવા માગી શકીએ. આ કિસ્સામાં $\Nat \to \Omega$ પ્રકારની વસ્તુઓને
એકસ્થાની સંબંધો અથવા~$\Nat$ના ઉપગણો તરીકે; $\Nat \to \Nat$ પ્રકારની
વસ્તુઓને પ્રાકૃતિક સંખ્યાઓથી પ્રાકૃતિક સંખ્યાઓમાં જતાં વિધેયો તરીકે;
અને $(\Nat \to \Nat) \to \Nat$ પ્રકારની વસ્તુઓને ``વિધેયકો'', એટલે
વિધેયોને સંખ્યાઓમાં મોકલતા ઉચ્ચ-પ્રકાર વિધેયો તરીકે વિચારી શકો.

દ્વિતીય-ક્રમ તર્કશાસ્ત્રની જેમ ઉચ્ચ-ક્રમ તર્કશાસ્ત્રને પણ બહુ-પ્રકાર
તર્કશાસ્ત્રના એક સ્વરૂપ તરીકે વિચારી શકાય, જેમાં આપણે વિચારવા માગીએ
એવા દરેક પ્રકારની વસ્તુ માટે એક પ્રકાર હોય. પરંતુ સામાન્ય રીતે
ઉચ્ચ-પ્રકાર તર્કશાસ્ત્રની વાક્યરચનાને પાયાથી જ વ્યાખ્યાયિત કરવી વધુ
સ્પષ્ટ છે. દાખલા તરીકે, સાન્ત પ્રકારોના ગણને નીચે મુજબ અનુમાનાત્મક રીતે
વ્યાખ્યાયિત કરી શકીએ:
\begin{enumerate}
\item $\Nat$ સાન્ત પ્રકાર છે.
\item જો $\sigma$ અને~$\tau$ સાન્ત પ્રકારો હોય, તો $\sigma
  \to \tau$ પણ સાન્ત પ્રકાર છે.
\item જો $\sigma$ અને~$\tau$ સાન્ત પ્રકારો હોય, તો $\sigma \times
  \tau$ પણ સાન્ત પ્રકાર છે.
\end{enumerate}
સહજ રીતે, $\Nat$ પ્રાકૃતિક સંખ્યાઓનો પ્રકાર સૂચવે છે,
$\sigma \to \tau$ એ~$\sigma$માંથી~$\tau$માં જતા વિધેયોનો પ્રકાર
સૂચવે છે, અને $\sigma \times \tau$ એવી વસ્તુઓના જોડાનો પ્રકાર સૂચવે
છે જેમાં એક વસ્તુ~$\sigma$માંથી અને એક~$\tau$માંથી આવે છે. ત્યાર પછી
પદોના ગણને નીચે મુજબ અનુમાનાત્મક રીતે વ્યાખ્યાયિત કરી શકીએ:
\begin{enumerate}
\item દરેક પ્રકાર~$\sigma$ માટે પ્રકાર~$\sigma$ના $x$, $y$, $z$, \dots ચલોનો
  જથ્થો હોય છે.
\item $\Obj 0$, પ્રકાર~$\Nat$નું પદ છે.
\item $\Obj S$ (ઉત્તરગામી), પ્રકાર $\Nat \to \Nat$નું પદ છે.
\item જો $s$, પ્રકાર~$\sigma$નું પદ હોય અને $t$, પ્રકાર
  $\Nat \to (\sigma \to \sigma)$નું પદ હોય, તો $\Obj{R}_{st}$,
  પ્રકાર $\Nat \to \sigma$નું પદ છે.
\item જો $s$, પ્રકાર $\tau \to \sigma$નું પદ હોય અને $t$,
  પ્રકાર~$\tau$નું પદ હોય, તો $s(t)$, પ્રકાર~$\sigma$નું પદ છે.
\item જો $s$, પ્રકાર~$\sigma$નું પદ હોય અને $x$, પ્રકાર~$\tau$નો ચલ
  હોય, તો $\lambd[x][s]$, પ્રકાર $\tau \to \sigma$નું પદ છે.
\item જો $s$, પ્રકાર~$\sigma$નું પદ હોય અને $t$, પ્રકાર~$\tau$નું પદ
  હોય, તો $\tuple{s, t}$, પ્રકાર $\sigma \times \tau$નું પદ છે.
\item જો $s$, પ્રકાર $\sigma \times \tau$નું પદ હોય, તો $p_1(s)$,
  પ્રકાર~$\sigma$નું અને $p_2(s)$, પ્રકાર~$\tau$નું પદ છે.
\end{enumerate}
સહજ રીતે, $\Obj{R}_{st}$ નીચે મુજબ આવર્તક રીતે વ્યાખ્યાયિત વિધેય
સૂચવે છે:
\begin{align*}
\Obj{R}_{st}(0) & = s \\
\Obj{R}_{st}(x+1) & = t(x, R_{st}(x)),
\end{align*}
$\tuple{s, t}$ એવો જોડો સૂચવે છે જેનો પ્રથમ ઘટક~$s$ અને બીજો ઘટક~$t$
છે, અને $p_1(s)$ તથા~$p_2(s)$,~$s$ના પ્રથમ અને બીજા ઘટકો
(``પ્રક્ષેપો'') સૂચવે છે. છેલ્લે, $\lambd[x][s]$ નીચેના વિધેય~$f$ને
સૂચવે છે:
\[
f(x) = s
\]
પ્રકાર~$\tau$ના કોઈ પણ~$x$ માટે; તેથી બાબત~(6) આપણને
ગુણધર્મ-ગણરચનાનું એક સ્વરૂપ આપે છે અને પદો વડે વિધેયો વ્યાખ્યાયિત કરવા
દે છે.
\sourcecorrection{OLBYD-003}{મૂળની
\texttt{higher-order-logic.tex}, પંક્તિ~87માં~$x$ને પ્રકાર~$\sigma$નો
કહ્યો છે, પરંતુ બાબત~(6)માં બાંધેલો ચલ~$x$ પ્રકાર~$\tau$નો છે અને
$\lambd[x][s]$નો પ્રકાર $\tau\to\sigma$ છે. અહીં પ્રકાર~$\tau$
પુનઃસ્થાપિત કર્યો છે.}
!!{formula}s સમાન પ્રકારનાં પદો વચ્ચેના !!{identity} વિધાનો
$\eq[s][t]$, સામાન્ય વિધાનાત્મક સંયોજકો અને ઉચ્ચ-પ્રકાર પરિમાણીકરણમાંથી
બને છે. પછી ઉપર વ્યાખ્યાયિત પદોને નિયંત્રિત કરતાં પાયાનાં સમીકરણો અને
પરિમાણકો તથા !!{identity} સાથેના તર્કશાસ્ત્રના સામાન્ય નિયમોને તંત્રની
સ્વયંસિદ્ધિઓ તરીકે લઈ શકાય છે.
\sourcecorrection{OLBYD-004}{મૂળની
\texttt{higher-order-logic.tex}, પંક્તિ~89માં શબ્દકોશ-ટોકન
\texttt{\string!\string!\string^\{formula\}s} તરીકે વિકૃત છે. બીજા બધા
પ્રકરણોની સુમેળવાળી ટોકન-રચના અનુસાર અહીં
\texttt{\string!\string!\{formula\}s} પુનઃસ્થાપિત
કર્યું છે.}

સાન્ત-પ્રકાર તંત્રમાં સત્યમૂલ્યોનો પ્રકાર~$\Omega$ ઉમેરવામાં આવે, તો
તેના ઉપયોગને નિયંત્રિત કરતી સ્વયંસિદ્ધિઓ પણ સામેલ કરવી પડે. હકીકતમાં,
ચતુરાઈથી કામ લઈએ તો જટિલ !!{formula}sને સંપૂર્ણપણે દૂર કરીને તેમની
જગ્યાએ પ્રકાર~$\Omega$નાં પદો મૂકી શકીએ!{} પછી સાબિતી-તંત્રને અનુરૂપ
રીતે બદલી શકાય. તેનું પરિણામ મૂળભૂત રીતે 1930ના દાયકામાં એલોન્ઝો
ચર્ચે રજૂ કરેલો \emph{પ્રકારોનો સરળ સિદ્ધાંત} છે.

દ્વિતીય-ક્રમ તર્કશાસ્ત્રની જેમ ઉચ્ચ-પ્રકાર અર્થવિજ્ઞાનનાં પણ જુદાં
સ્વરૂપો વાપરી શકાય છે. પૂર્ણ સ્વરૂપમાં, પ્રકાર $\sigma \to \tau$ના ચલો
પ્રકાર~$\sigma$ની વસ્તુઓમાંથી પ્રકાર~$\tau$ની વસ્તુઓમાં જતા \emph{બધા}
વિધેયોના ગણ પર વ્યાપે છે. ધાર્યા મુજબ આ અર્થવિજ્ઞાન એટલું શક્તિશાળી છે
કે તે કોઈ પૂર્ણ, અસરકારક !!{derivation} તંત્રને સ્વીકારતું નથી. પરંતુ
દુર્બળ અર્થવિજ્ઞાન પણ વિચારી શકાય, જેમાં !!a{structure} દરેક
પ્રકાર~$\tau$ માટે ઘટકોનો ગણ~$T_\tau$ તથા પ્રયોજન, પ્રક્ષેપ
વગેરે માટેની યોગ્ય ક્રિયાઓ ધરાવે છે. વિગતો યોગ્ય રીતે ઘડીએ, તો ઉપર
વર્ણવેલા પ્રકારનાં !!{derivation} તંત્રો માટે પૂર્ણતા પ્રમેયો મેળવી
શકાય છે.

ઉચ્ચ-પ્રકાર તર્કશાસ્ત્ર આકર્ષક છે, કારણ કે તે એવું માળખું આપે છે જેમાં
ગણિતનો મોટો ભાગ કુદરતી રીતે સમાવી શકીએ: $\Nat$થી શરૂ કરીને વાસ્તવિક
સંખ્યાઓ, સતત વિધેયો વગેરે વ્યાખ્યાયિત કરી શકીએ. સ્ફુરણાવાદી
તર્કશાસ્ત્રના સંદર્ભમાં પણ તે ખાસ આકર્ષક છે, કારણ કે પ્રકારોનાં સ્પષ્ટ
``રચનાત્મક'' અર્થઘટનો છે. હકીકતમાં, ઉચ્ચ-પ્રકાર અર્થવિજ્ઞાનનાં
રચનાત્મક સ્વરૂપો (પ્રશિષ્ટ તર્કશાસ્ત્રને બદલે સ્ફુરણાવાદી તર્કશાસ્ત્ર
પર આધારિત) વિકસાવી શકાય છે; તે આ રચનાત્મક અર્થઘટનોને સુંદર રીતે સ્પષ્ટ
કરે છે અને ઘણી રીતે પોતાના પ્રશિષ્ટ પ્રતિરૂપો કરતાં વધુ રસપ્રદ છે.

\end{document}
